\documentclass{article}



\usepackage[utf8]{inputenc}
\usepackage[english]{babel}
\usepackage{xcolor}
\usepackage{amssymb}
\usepackage{lineno}
\usepackage{authblk}
\usepackage{natbib}
\usepackage[margin=2.5cm]{geometry}
\usepackage{setspace}
\usepackage{chemformula}

% Language setting
% Replace `english' with e.g. `spanish' to change the document language
\usepackage[english]{babel}

% Set page size and margins
% Replace `letterpaper' with `a4paper' for UK/EU standard size
%\usepackage[letterpaper,top=2cm,bottom=2cm,left=3cm,right=3cm,marginparwidth=1.75cm]{geometry}

% Useful packages
\usepackage{amsmath}
\usepackage{graphicx}
\usepackage[colorlinks=true, allcolors=blue]{hyperref}
\usepackage{natbib}


\title{Your Paper}
\author{You}

\begin{document}
\maketitle

\begin{abstract}
Your abstract.
\end{abstract}

\section{Introduction}



%Comparison with OCO-2 and CarboScope inversions also indicates a substantial improvement of the global mean sea- sonal cycle of NEE (Fig. 5)


%The assess- ment of the net land–atmosphere exchange from DGVMs and atmospheric inversions also shows substantial discrep- ancy, particularly for the estimate of the net land flux over the northern extra-tropics. This discrepancy highlights the dif- ficulty of quantifying complex processes (CO2 fertilization, nitrogen deposition and fertilizers, climate change and vari- ability, land management, etc.) that collectively determine the net land CO2 flux. Resolving the differences in the North- ern Hemisphere land sink will require the consideration and inclusion of larger volumes of observations. Eddy covariance data consisted of 294 sites from around the world though skewed towards a higher representation of tem- perate forests from North America and Europe.For the generation of global flux maps we used hourly meteo- rological data from ERA5 global reanalysis products at 0.25° . measurements of the MODerate Imaging Spectroradiometer (MODIS) of surface reflectance and land surface temperature from collection 006 at daily resolution. Machine learning All X-BASE products are based on gradient-boosted regres- sion trees using the XGBoost library (Chen and Guestrin, 2016). X\cite{Nelson2024}



%The X-BASE product estimates the global terrestrial NEE to be −5.75+-0.33 PgC yr-1 (2001–2020), with hotspots of strong CO2 flux into the ecosystems in the tropical re- gions and temperate regions of North America and Eu- rope (Fig. 4). In contrast to both RS and RS+METEO, India and some regions in central Sahel show prominent patterns of a mean CO2 flux from the ecosystems to the atmosphere in X-BASE, corresponding mostly to crop- designated areas (Fig. C2). 

%However, X-BASE global ter- restrial NEE (−5.63 PgC yr-1) agrees well with the inver- sion estimates of OCO-2 (−4.12 PgC yr-1) and CarboScope (−3.46 PgC yr-1) over the common period (2015–2020). Here, GFED 4.1 (Randerson et al., 2017) emission estimates have been used to correct the inversion estimates for fire emissions. Comparison with OCO-2 and CarboScope inversions also indicates a substantial improvement of the global mean sea- sonal cycle of NEE (Fig. 5) in X-BASE compared to RS and RS+METEO. The systematic bias present in RS and RS+METEO has essentially disappeared in X-BASE. The shape, and in particular the amplitude, of the global NEE seasonal cycle of X-BASE is more consistent with the in- versions. The larger and more realistic seasonal cycle am- plitude of global NEE in X-BASE originates primarily from improved and increased amplitudes in boreal regions.






% -- Prioritization of areas.
% have implications for soil management for mitigation of global warming. through carbon sequestration. First, the results indic- ate a range of possible rates of global SOC stock change (figure 9; see figure S8 for zonal results). Second, the spatial results showed which regions are likely to respond more sensitively, specific- ally, to land-use change and land management (e.g. figures S5–S7). Third, the simulation results provide supple mentary evidence that prevention of deforestation or reforestation can enhance SOC stocks.  \citep{Ito2019}

% The global nature of carbon sequestration supports the implementation of diverse, ecosystem-based approaches to climate change mitigation across different latitudes.




% Forests the same

% Forests, especially tropical forests, play a vital role in the global carbon cycle, holding back a significant amount of atmospheric warming. Boreal forests face threats from increased wildfires and logging, which release significant amounts of stored carbon.

% The higher CS values observed in forests (Fig.~\ref{fig:CMIP6_NEP_map}a) underscore the important role of these ecosystems in the global carbon cycle. This means that carbon inputs have been larger than carbon outputs during this time horizon, highlighting the capacity of terrestrial ecosystems to contribute to mitigating anthropogenic carbon emissions. 
% Forest conservation priority: The high carbon sequestration rates in forests, particularly in the tropics, emphasize the urgent need for conservation and sustainable management of these ecosystems.

%  Areas outside the tropics also show positive sequestration, suggesting that reforestation and afforestation efforts in temperate and boreal regions can contribute meaningfully to carbon mitigation.




% Estimates that take cost into account suggest that reforestation could increase the land sink by up to 8–13.8 Pg of CO2 equivalents (CO2e) per year (2.3–3.7 PgC yr-1) globally from 2020 to 2050260. This value could even be underestimated if CO2 and climate effects continue to enhance the land sink (for example, ref. 22). However, the value could be overestimated if the impacts of disturbances or cost become large. Reforestation, agroforestry and reducing deforestation together have the potential to mitigate up to 20 PgCO2e yr-1 (Fig. 7a)260 . The main policy tool aimed at reforestation and reduced deforestation is carbon credits22 , which assign monetary value to carbon captured from, or emissions prevented to, the atmosphere. The general consensus is that protecting old-growth forests is optimal for carbon storage because existing forests, even when old, continue to act as net land carbon sinks and represent large stores of avoided emissions78,81 . However, there are potential downsides of NbCS projects, despite their potential to store carbon. First, the carbon sink in planted forests could be less persistent than primary forests, owing to wind throw events and fire regimes212,262–265 . Second, afforestation, reforestation and forest protection programmes have been criticized for worsening inequality in the Global South266 , and carbon credit programmes are liable to abet greenwashing. Gains from NbCS projects are also modest relative to current anthropogenic emissions of fossil carbon. On the whole, however, NbCS projects, when appropriately implemented at scale, are expected to increase the land carbon sink and slow the growth of atmospheric CO2 \cite{Ruehr2023}


%Key Factors Influencing Carbon Fluxes: Soil moisture is a better predictor of detrended interannual variability of Net Biome Production (NBP) than temperature in most regions, except for the Amazon. Climate change impacts can negatively affect land, but the efficiency of land in removing CO2 has remained broadly constant. Land-use change emissions have a large uncertainty. Deforestation prevention and reforestation can enhance SOC stocks. Forest conservation, especially in the tropics, is crucial due to high carbon sequestration rates. There is a tension between forests and agricultural expansion, particularly in the Amazon.

%Wildfire—although a natural element in boreal forests—represents one of the greatest threats to boreal forest carbon. With increased temperatures, rising more than twice as fast in boreal forests compared to lower latitudes, and more frequent and long-lasting droughts, boreal forests are now experiencing more frequent and intense wildfires. The hotter and more often a stand of boreal forest catches fire, the deeper into the soil carbon pool the fire will burn, sending centuries-old carbon up in smoke in an instant. Logging of high-carbon primary forests is also a big issue in the boreal.

%the emissions from permanent deforestation – the largest of our component fluxes – could be halted (largely) without compromising carbon uptake by forests, contributing substantially to emissions reduction. By con- trast, reducing wood harvesting would have limited poten- tial to reduce emissions as it would be associated with less forest regrowth; removals and emissions cannot be decou- pled here on long timescales. Although the budget imbalance is near zero for the recent decades, this could be due to compensation of errors. W \cite{Friedlingstein2023}


%Despite these regional negative effects of climate change on SLAND, the efficiency of land to remove anthropogenic CO2 emissions has remained broadly constant over the last 6 decades, with a land-borne fraction (SLAND/(EFOS + ELUC)) of around 30\% (Fig. 9b). \cite{Friedlingstein2023}

%Global fossil CO2 emissions, EFOS (including the cement carbonation sink), have increased every decade from an av- erage of 3.0+-0.2 GtC yr-1 for the decade of the 1960s to an average of 9.6+-0.5 GtC yr-1 during 2013–2022 (Table 7, Figs. 2 and 5). The growth rate of these emissions decreased between the 1960s and the 1990s, from 4.3\%yr-1 in the 1960s (1960–1969), 3.2\%yr-1 in the 1970s (1970-1979), 1.6\%yr-1 in the 1980s (1980–1989) to 1.0\%yr-1 in the 1990s (1990–1999). After this period, the growth rate be- gan increasing again in the 2000s at an average growth rate of 2.8\%yr-1, decreasing to 0.5\%yr-1 for the last decade (2013–2022). The cumulative land sink is almost equal to the cumulative land-use emissions (220+-70 Gt C), making the global land nearly neutral over the whole 1850– 2022 period\cite{Friedlingstein2023}

%Models show that after the CO2 concentration and air temperature peaks, land and ocean are decreasing carbon sinks from the 2,040s and become sources for a limited time in the 22nd century.The decrease in the carbon uptake precedes the CO2 concentration peak. The early peak of the land uptake occurs due to a larger increase in ecosystem respiration than the increase in gross primary production, as well as due to a concomitant increase in land-use change emissions primarily attributed to the wide implementation of biofuel croplands. \citet{Melnikova2021}

%Overall, the models suggest that the global and NHL terrestrial ecosystems will be consistent carbon sinks in the future, and the magnitude of the carbon sinks is projected to be larger under scenarios with higher radiative forcing. \citep{Qiu2023}

%Increasing atmospheric CO2 concentrations have increased the land C sink by enhancing GPP more than ecosystem respiration. The reduction in S deposition rates is severely altering the C balance by enhancing photosynthesis and ecosystem respiration but in a decoupled manner. It is far from certain whether terrestrial ecosystems will continue to respond positively to increasing CO2, will saturate, or will eventually reach a tipping point beyond which respiration and the release of greenhouse gases exceed production.\cite{Fernandez-Martinez2017}



%Carbon management can be oriented to maximize CS and CBS, which can be achieved by managing both rates of carbon input and process rates in ecosystems. We believe the use of these metrics can help to better deal with current discussions about the role of ecosystems in mitigating climate change, provide better estimates of avoided or human-induced warming, and have the potential to be included in accounting methods for climate policy.

% --General



% --- Time scale
% The balance of carbon assimilation and respiration is critical: The higher carbon sequestration in the tropics suggests that high tropical assimilation outweighs low respiration in extratropical regions when only climate is considered. Understanding the long-term balance between carbon assimilation (photosynthesis) and respiration (release of CO2) is vital for predicting future climate scenarios. Factors such as increased atmospheric CO2, higher nitrogen deposition, and climate change itself can affect this balance, with both positive and negative feedback loops.

% Time: Changes in vegetation carbon residence times can cause major shifts in the distribution of carbon between pools, overall fluxes, and the time constants of terrestrial carbon transitions, with consequences for the land carbon balance and the associated state of ecosystems.

% The loss of carbon from disturbance may not be released to the atmosphere in the year of loss but delayed for years (e.g., timber products) or transferred to other pools (e.g., live to dead) (43). \cite{Xu2021}

% The persistent unexplained variability in the carbon budget imbalance limits our ability to verify reported emissions (Peters et al., 2017) and suggests we do not yet have a com- plete understanding of the underlying carbon cycle dynam- ics on annual to decadal timescales The long-standing sparse data coverage of fCO2 observations in the South- ern Hemisphere compared to the Northern Hemisphere (e.g. Takahashi et al., 2009) continues to exist (Bakker et al., 2016, 2023, Fig. S1) and to lead to substantially higher uncertainty in the SOCEAN estimate for the Southern Hemisphere (Wat- son et al., 2020; Gloege et al., 2021; Hauck et al., 2023).\cite{Friedlingstein2023}




% -- Differences among products
















% En bosques donde ha habido un uso intensive como en Euroasia y Norteaméroica, los efectos se acumulan, hay como una "deuda" de carbón. Mientras en bosques más protegidos como en el Congo, el suroeste de asua y el Amazonas, hay cierta "ganancia".

% NEP: Factor climática. Efecto de las variaciones en el clima.
% NBP: efecto del uso del suelo y cambio del suelo. En algunos modelos se considers el fuego.

%Net Ecosystem Production denotes the net accumulation of organic matter or carbon by an ecosystem; NEP is the difference between the rate of production of living organic matter (NPP) and the decomposition rate of dead organic matter (heterotrophic respiration, RH). Heterotrophic respiration includes losses by herbivory and the decomposition of organic debris by soil biota. Global NEP is estimated to about 10 Gt C yr-1. NEP can be measured in two ways: One is to measure changes in carbon stocks in vegetation and soil; the other is to integrate the fluxes of CO2 into and out of the vegetation (the net ecosystem exchange, NEE) with instrumentation placed above (Aubinet et al., 2000). The precision of both of these methods is improving.

%Net Biome Production denotes the net production of organic matter in a region containing a range of ecosystems (a biome) and includes, in addition to heterotrophic respiration, other processes leading to loss of living and dead organic matter (harvest, forest clearance, and fire, etc.) (Schulze and Heimann, 1998). NBP is appropriate for the net carbon balance of large areas (100-1000 km2) and longer periods of time (several years and longer). In the past, NBP has been considered to be close to zero (Figure 1-2). Compared to the total fluxes between atmosphere and biosphere, global NBP is comparatively small; NBP for the decade 1989-1998 has been estimated to be 0.7 +- 1.0 Gt C yr-1 (Table 1-2)-about 1 percent of NPP and about 10 percent of NEP.

%NEE=NEP. These terms are used somewhat interchangeably, with NEE used more often to refer to these fluxes when they are addressed from a measurement of gas exchange rates using atmospheric measurements over time scales of hours, whereas NEP is more often used to refer to the same processes if measurements are based on ecosystem-carbon stock changes, typically measured over a minimal period of one year. However, these differences in usage are not firmly embedded in formal definitions.

\citet{Arora2020}
Here, we compare carbon– concentration and carbon–climate feedbacks from models participating in the C4MIP (Jones et al., 2016) contribution to the sixth phase of CMIP (CMIP6; Eyring et al., 2016). To maintain continuity and consistency, feedback parameters are derived from the 1pctCO2 experiments as

This paper briefly describes each model 

Model intercomparison projects offer several benefits includ- ing the calculation of the model mean response. 

The carbon uptake over land and ocean, in response to increasing atmospheric CO2 concentration, is well known to be dominated by the positive contribution from the carbon–concentration feedback (Gregory et al., 2009; Arora et al., 2013). The strength of this feedback is of comparable magnitudes over land (mean+-standard deviation=0.97+-0.40 PgCppm-1) and ocean (0.79+-0.07 PgCppm-1) although the feedback is much more uncertain over land as indicated by the standard deviation across the 11 models considered here. 

since the carbon– concentration and carbon–climate feedbacks are determined entirely by biological process, which are much less under- stood, the resulting uncertainty is much higher across the land models than across the ocean models.

Over land, the diverse re- sponse of models is found to be primarily due to the wide range of the strength of the CO2 fertilization effect, the frac- tion of GPP that is converted to NPP, and the residence times of carbon in the live (vegetation) and dead (litter plus soil) carbon pools across models. There

\citet{Gier2024}
Improvements in the representation of the land carbon cycle in Earth system models participating in the Coupled Model Intercomparison Project Phase 6 (CMIP6) include interactive treatment of both the carbon and nitrogen cycles, im- proved photosynthesis, and soil hydrology.  To assess the impact of these model developments on aspects of the global carbon cycle, the Earth System Model Evaluation Tool (ESMValTool) is expanded to compare CO2 concentration and emission driven historical simulations from CMIP5 and CMIP6 to observational data sets.

Overestimations of photosynthesis (gross primary productivity (GPP)) in CMIP5 were largely resolved in CMIP6 for participating models with an interactive nitrogen cycle, but remaining for models without one. This points to the importance of including nutrient limitation

In ESMs, global mean land carbon uptake (net biome productivity (NBP)) is well reproduced in the CMIP5 and CMIP6 multi-model means. However, \textcolor{teal}{this is the result of an underestimation of NBP in the northern hemisphere, which is compensated by an overestimation in the southern hemisphere and the tropics.} While vegetation carbon content is slightly better represented in CMIP6, the inter-model range of soil carbon content remains the same between CMIP5 and CMIP6.

This improvement is not as significant in the multi-model means due to more new models in CMIP6, especially those using older versions of the Community Land Model (CLM). Emission driven simulations perform just as well as concentration driven models despite the added process-realism.

Regridded to 2deg x 2deg. They did not say the method. They only regridded to calculate the mean. 

\citet{Melnikova2021}

Models show that after the CO2 concentration and air temperature peaks, land and ocean are decreasing carbon sinks from the 2,040s and become sources for a limited time in the 22nd century.The decrease in the carbon uptake
precedes the CO2 concentration peak. The early peak of the land uptake occurs due to a larger increase in ecosystem respiration than the increase in gross primary production, as well as due to a concomitant increase in land-use change emissions primarily attributed to the wide implementation of biofuel croplands.

All ESMs estimate lower land carbon uptake in Russia, South Asia, and Southeast Asia than the correspond- ing inventory-based estimates in RECCAP (Figure 1c). The inversions estimate a higher land sink in the north than ESMs and a carbon source in the tropics that is not present in the ESMs. MIROC-ES2L estimates a slightly higher net biome production (NBP) sink in Northern Hemisphere than other models (that is still lower than in the inversions). UKESM1-0-LL estimates slightly high carbon uptake in tropical Amazonia and Africa, which Sellar et al. (2019) attributed to a small overestimation of tree fraction in savanna biomes compared to observational data sets.

RECCAP. The difference between NPP and RH is higher in RECCAP estimates than in ESMs in all regions. Thus, these component fluxes cannot explain the discrepancies between ESM and RECCAP estimates, suggesting there could still be missing processes in these state-of-the-art models (Ciais et al., 2020).
The CMIP6 ESMs show higher historical land carbon uptake than GCB2019 and lower uptake than esti- mates by Ciais et al. (2020). This is probably due to the differences in terms included in the land sink by two evaluation approaches–anthropogenic sink in GCB2019 and anthropogenic and natural sinks in the study by Ciais et al. (2020). However, the use of both NBP anomaly from piControl (Figure 3c) and NBP absolute value (Figure S4d) lead to very similar results. Whether the natural or anthropogenic terms of land uptake in certain regions were underestimated requires further analysis and more attention in future studies.

\citet{Hu2022}
In this work, we studied the model performance of 22 ESMs participating in the fifth and sixth phases of the Coupled Model Intercomparison Project (CMIP5 and CMIP6) by comparing historical GPP and NPP simulations with satellite data fromMODIS and further evaluating potential model improvement fromCMIP5 to CMIP6.
InCMIP6, the average global total GPP andNPP estimatedby the 22 ESMsare 16\%and 13\%higher thanMODIS data, respectively. The multi-model ensembles (MME) of the 22 ESMs can fairly reproduce the spatial distribution, zonal distribution and seasonal varia- tions of both GPP and NPP from MODIS. They perform much better in simulating GPP and NPP for grasslands, wet- lands, croplands and other biomes than forests. However, there are noticeable differences among individual ESM simulations in terms of overall fluxes, temporal and spatial flux distributions, and fluxes by biome and region. The MME consistently outperforms all individual models in nearly every respect.

The carbon flux data of both ESMs were regridded to the same resolu-
tion of 0.5° × 0.5° by using bilinear interpolation. From the begining.

\textcolor{magenta}{To explain the uncertainties between models, we also introduced model spread (MS) (Huang et al., 2018). Eq. 1}

The overall simulation effectiveness of GPP is better than that of NPP. The MME consistently outperforms individual models. Each model has its own advantages and disadvantages in different aspects of simulation, largely attributable to the different land surface schemes, resolutions and parameterizations. The ESMs perform differently by biome type, with the worst performance in forest ecosystems.

\citep{Ito2019}
In this study, we analyzed the outputs of simulations ofSOC stock by Earth system models (ESMs), most ofwhich are participants in the Land-Use Model Intercomparison Project. the ESMs simulated SOC distribution patterns and their changes during historical (1850–2014) and future (2015–2100) periods. Total SOC stock increased in many simulations over the historical period (30 ± 67 Pg C) and under future climate and land-use conditions (48 ± 32 Pg C for ssp126 and 49 ± 58 Pg C for ssp370). Although there were considerable inter-model differences, the rates are still remarkable in terms of their potential for mitigation of global warming. The disparate results among ESMs imply that key parameters that control processes such as SOC residence time need to be better constrained and that more comprehensive representation of land management impacts on soils remain critical for understanding the long-term potential of soils to sequester carbon.

On average, cLitter + cCwd accounted for 11.9% of SOCtotal; the global cLitter + cCwd fraction ranged from 1.7% to 27.8% among the ESMs.
Most models simulated high carbon stocks in the northern high latitudes and low carbon stocks in the middle to low latitudes (figure 2; see figure S2 for cSoil maps), con- sistent with findings by Todd-Brown et al (2013) and Carvalhais et al (2014) for CMIP5 models. This sim- ulated latitudinal gradient reflects variations in the decomposition rate and is consistent with large-scale patterns seen in observational soil carbon datasets.

Global mean residence time of SOC total in the
1850s was calculated as 36.8 ± 20.5 years and ranged from 11.4 years in GFDL-ESM4 to 84.1 years in CMCC-CM2-SR5 (figure 1). In most models, the res- idence time fell within the observational range of 18.5–45.8 years (Raich and Schlesinger 1
The difference in mean residence time explains about 88% of the variation of global SOC stock among the models (figure S3). Previous studies have pointed out that the turnover rate is an important metric ofecosystemcar- bon pools (e.g. Bonan et al 2013, Todd-Brown et al 2013, Carvalhais et al 2014, Friend et al 2014). For

The highest accumulation rate was simulated by BCC-CSM2-MR: 1.1 Pg C yr-1 or 0.66\%0 yr-1. Historical changes in SOC stocks occurred het-
erogeneously over the land surface and differed among ESMs (figure S5). Three models (ACCESS- ESM1-5, BCC-CSM3-MR, and CNRM-ESM2-1) showed extensive increases, whereas two models (EC-Earth3-Veg and GFDL-ESM4) that used gross land-use transition showed decreases in most land areas. The

In Europe, for example, the majority of models simulated a net increase ofSOC,whereasCanESM5andGFDL-ESM4 simulated regional decreases ofSOC stocks except in a few countries. InNorth America (MidwesternUnited States), the simulated historical change was obvi- ously inconsistent among the models: clear increases were simulated by six models, clear decreases by five models, and mixed patterns by four models. Similar inconsistencies were simulated in other cultivated areas such as in central South America and East Asia (China). I

Although the model results presented here appear inconsistent, nonetheless have implications for soil management for mitigation of global warming. through carbon sequestration. First, the results indic- ate a range of possible rates of global SOC stock change (figure 9; see figure S8 for zonal results). Second, the spatial results showed which regions
are likely to respond more sensitively, specific- ally, to land-use change and land management (e.g. figures S5–S7). Third, the simulation results provide supple-
mentary evidence that prevention of deforestation or reforestation can enhance SOC stocks.

Overall, the results indicate that total SOC is likely to increase, although at a modest rate (about 0.4\%0 yr-1). 

In addition, these modeling efforts emphasize the importance ofmonitoring and verification ofSOC at global scale (Smith et al 2020).

\citep{Koch2021}
Recent analyses of long-term (1985–2014) forest inventory plots across the tropics show that structurally intact tropical forest are a large carbon sink, but that this sink has saturated and is projected to be in long-term decline. Here we compare these results with estimates from the two latest generations of Earth System Models, Climate Modelling Intercomparison Project 5 (CMIP5) (19 models) and CMIP6 (17 models). While CMIP5 and CMIP6 are of similar skill, they do not reproduce the observed 1985–2014 carbon dynamics. CMIP multimodel-means reproduce the response of carbon gains from tree growth to environmental drivers, but the modeling of carbon losses from tree mortality does not correspond well to the inventory data. The model-observation differences primarily result from the treatment of mortality in models.
This resulted in a 66% smaller total natural pan-tropical sink in the newest generation state-of-the-art CMIP6 models (0.34 ± 1.15 Pg C yr−1), as compared to the observed pan-tropical sink from 1985-2014 (1.02 ± 0.40 Pg C yr−1). Furthermore, the models show a moderate increase in carbon uptake in both the African and Amazonian tropics over the 1985–2014 window, whereas there is a clear decline in the sink in Amazonia. Future projections to 2040 show an increase in the sink to 2040 using the CMIP6 models under all four diverse future scenarios analyzed. This contrasts starkly with the observed strong past decline in the Amazonian carbon sink and the statistically derived predicted future decline of the natural carbon sink on both continents.
Critically, these mismatches can be explained because models do not correctly balance CO2 fertilisation effects and respiration costs, while they also fail to capture trends in mortality, likely stemming from the dif- ficulties of modeling the complex lag effects between past carbon gains and carbon losses. This

\citep{Padron2022}
Nevertheless, the latest cumulative land carbon sink projections from 11 Earth system models participating in the Coupled Model Intercomparison Project phase 6 (CMIP6) show an intermodel spread of 150 PgC (i.e., appr15 years of current anthropogenic emissions) for a policy-relevant sce- nario, with mean global warming by the end of the century below 2 degC relative to preindustrial conditions. We
Using multiple linear regression and a resampling technique, we quantify the individual contributions of these five drivers for explaining the cumulative NBP anomaly of each model relative to the multi-model mean. We find that the intermodel variability of the contributions of each driver relative to the total NBP intermodel variability is 52.4% for the sensitiv- ity to temperature, 44.2% for the sensitivity to soil mois- ture, 44% for the sensitivity to CO2, 26.2% for the average temperature, and 21.9% for the average soil moisture. Fur- thermore, the sensitivities of NBP to temperature and soil moisture, particularly at tropical regions, contribute to ex- plain 34% to 65% of the cumulative NBP deviations from the ensemble mean of the two models.

Even under this sce- nario with a relative low concentration of greenhouse gases, there is an intermodel spread of approximately 150 PgC in cumulative NBP from 2015 to 2100 – equivalent to 15 years of current anthropogenic emissions – which translates into a 40% uncertainty in the carbon budget remaining to stabi- lize global temperature below the chosen threshold. We also show that even when two models project a similar global land carbon sink, there can be large and compensating regional differences.

In addition, we find that the detrended interannual variability of projected NBP is better explained by soil moisture than temperature in most models and across most regions. A notable excep- tion is the core of the Amazon, where temperature is more important than soil moisture to explain the interannual vari- ability of NBP in the models.

This finding provides explicit evidence that improving the representation of the local-scale sensitivity of NBP to interannual climate variability has the potential to re- duce uncertainty in long-term projections of the global land carbon sink. For example, a high regional contribution of sSM and/or sT to a model’s land sink anomaly indicates the need to evaluate and improve potentially related processes such as water stress on photosynthesis, the effect of temperature and moisture on soil carbon loss due to microbial activity, and the occurrence and magnitude of fire emissions.

\citep{Qiu2023}
Overall, the models suggest that the global and NHL terrestrial ecosystems will be consistent carbon sinks in the future, and the magnitude of the carbon sinks is projected to be larger under scenarios with higher radiative forcing. By the end of this century, the models on average estimate the NHL net ecosystem productivity (NEP) as 0.54+-0.77, 1.01+-0.98, 0.97+-1.62, and 1.05±1.83 PgC yr-1 under SSP126, SSP245, SSP370, and SSP585, respectively. The uncertainties are not substantially reduced compared with earlier results, e.g., the Coupled Climate–Carbon Cycle Model Intercomparison Project (C4MIP). Although NHLs contribute a small fraction of the global carbon sink (app13 %), the relative uncertainties in NHL NEP are much larger than the global level. Our

Based on the average of the CMIP6 models, our analysis showed that global and NHL ecosystems were and would continue to be carbon sinks, although large uncertain- ties were found for the size and trends of the carbon sinks among different CMIP6 models, which are not obviously attenuated compared with previous model intercomparison project results. Although the warming levels and sensitiv- ity of ecosystems to the warming temperature are higher in the NHLs, the contribution of NHLs to the global NEP in- crease is small, however with larger relative uncertainties.

\citep{Shi2024}
. Here, we assessed the reliability of Earth system model (ESM) predictions of soil C change using the Coupled Model Intercomparison Project phases 5 and 6 (CMIP5 and CMIP6). ESMs predicted global soil C gains under the high emission scenario, with soils taking up 43.9 Pg (95% CI: 9.2–78.5 Pg) C on average during the 21st century. The variation in global soil C change declined significantly from CMIP5 (with average of 48.4 Pg [95% CI: 2.0–94.9 Pg] C) to CMIP6 models (with average of 39.3 Pg [95% CI: 23.9–54.7 Pg] C). For some models, a small C increase in all biomes contributed to this convergence. For other models, offsetting responses between cold and warm biomes contributed to convergence. Although soil C predictions appeared to converge in CMIP6, the dominant processes driving soil C change at global or biome scales differed among models and in many cases between earlier and later versions of the same model. Random Forest models, for soil carbon dynamics, accounted for more than 63% variation of the global soil C change predicted by CMIP5 ESMs, but only 36% for CMIP6 models. Although most CMIP6 models apparently agree on increased soil C storage during the 21st century, this consensus obscures substantial model disagreement on the mechanisms underlying soil C response, calling into question the reliability of model predictions.

Even though ESM predictions now cluster more tightly around the mean soil C gain, they all changed something different. These underlying issues suggest that the models have not in fact reached a consensus and that predictive confidence is low for soil C feedbacks to climate change. Moreover, uncertainty has not yet peaked because the models are still missing key processes.

\citep{Spafford2021}
Despite these common experimental protocols, a variety of terrestrial biogeochemical cycle validation approaches were adopted by CMIP6 participants, leading to ambiguous model perfor- mance assessment and uncertainty attribution across ESMs. In this review we summarize current methods of terrestrial biogeochemical cycle validation utilized by CMIP6 partici- pants and concurrent community model comparison studies. To address this inconsistency and alleviate the immense responsibilities of participants, we recommend the designation of a standard validation proto- col for CMIP participants, which is consolidated in an open- source software (such as the Earth System Model Evalua- tion Tool version 2 (ESMValToolv2.0) or the International Land Model Benchmarking version 2.1 (ILAMBv2.1)). This protocol should utilize a comprehensive suite of certainty- weighted observational reference datasets, targeted model performance metrics, and comparisons across a range of spatiotemporal dimensions.

\citep{Tokarska2020}

Here we quantify the impact of internal variability on the carbon budgets consistent with the Paris Agreement derived using either approach, and on the TCRE diagnosed from individual models. Our results show that internal variability contributes approximately +-0.09 ◦C to the overall uncertainty range of the human-induced warming to-date, leading to a spread in the remaining carbon budgets as large as +-50 PgC

We do not question the validity of a carbon budget approach in determining mitigation requirements. However, due to intrinsic uncertainty arising from internal variability, it may only be possible to determine the exact year when a budget is exceeded in hindsight, highlighting the importance of a precautionary approach

However, our ability to estimate the anthropogenic component of the observed warming is hampered by the presence of internal variability in the climate system. Unlike the anthropogenic warming component simulated by climate models (which, in principle, could be free from internal climate variability given a sufficiently large ensemble of model simulations), we only have one observed record that is modulated by natural variability and further subject to considerable measurement uncertainty. Thus, the presence of natural variability in the real world intrinsically limits our ability to determine the anthropogenic component of ‘where we are’ in any single year


\citep{Varney2022}
\textcolor{magenta}{Read it complete}
In this study CMIP6 ESMs are evaluated against empirical datasets to assess the ability of each model to simulate soil carbon and related controls: net primary productivity (NPP) and soil carbon turnover time (taus). Comparing CMIP6 with the previous generation of models (CMIP5), a lack of consistency in modelled soil carbon remains, particularly the underestimation of northern high-latitude soil carbon stocks. There is a robust improvement in the simulation ofNPP in CMIP6 com-
pared with CMIP5; however, an unrealistically high correlation with soil carbon stocks remains, suggesting the potential for an overestimation of the long-term terrestrial carbon sink. Additionally, the same improvements are not seen in
the simulation of taus. These results suggest that much of the uncertainty associated with modelled soil carbon stocks can be attributed to the simulation of below-ground processes,and greater emphasis is required on improving the representation of below-ground soil processes in future developments of models.


The spatial patterns of soil carbon in CMIP6 models appear to be more in agreement with each other than they were in CMIP5 and are more consistent with ob-
servations in the mid-latitudes. However, soil carbon is still heavily underestimated in high northern latitudes (with the exception of the two CMIP6 models that represent deep soil carbon).

There is good evidence that spatial patterns of contemporary NPP are better simulated in CMIP6 than in CMIP5 generation models when compared with satellite-derived estimates.

spatial patterns of taus continue to be poorly represented in CMIP6 models, in comparison with estimates derived from observational datasets of soil carbon
and heterotrophic respiration (Rh).

soil carbon simulations in both the CMIP5 and CMIP6 ESM generations seem to be spuriously highly correlated with NPP, which may make soil carbon in these models over-responsive to future projected changes in NPP.

\citep{Varney2023}
This study evaluates projections of soil carbon during the 21st century in Coupled Model Intercomparison Project Phase 6 (CMIP6) Earth system models (ESMs) under a range of atmospheric composition scenarios. In general, we find a reduced spread of changes in global soil carbon (deltaCs) in CMIP6 compared to the previous CMIP5 model generation. However, similar reductions were not seen in the derived contributions to deltaCs due to both increases in plant net primary productivity (NPP,named deltaCs,NPP) and reductions in the effective soil carbon turnover time (taus, named deltaCs,tau ). Instead, we find a strong relationship across the CMIP6 models between these NPP and taus components of deltaCs, with more positive values of deltaCs,NPP being correlated with more negative values of deltaCs,tau . We
show that the concept of “false priming” is likely to be contributing to this emergent relationship, which leads to a decrease in the effective soil carbon turnover time as a direct result of NPP increase and occurs when the rate of increase in NPP is relatively fast compared to the slower timescales of a multi-pool soil carbon model.

The derived linear terms which contribute to net soil carbon change, the response of soil carbon due to changes in NPP (deltaCs,NPP), and the response due to changes in taus (deltaCs,tau ) are found to have a strong relationship in CMIP6, with a more significant correlation than what was seen in CMIP5. This correlation is likely to be a cause of the reduction in the deltaCs projection spread
across the CMIP6 ensemble.

False priming was found to likely be contributing to the apparent emergent relationship between deltaCs,NPP and deltaCs,tau in CMIP6 ESMs, which describes a transient re-
duction in effective soil carbon turnover time due to increased input of carbon to the soil. The net effect of false priming is a coupling affect between deltaNPP and deltataus and results in a reduced range of future deltaCs predictions in CMIP6.

Our study highlights the significant role that false priming can play under transient changes in atmospheric CO2 and climate.We advise caution in the interpretation of changes in the effective soil carbon turnover time in terms of climate effects alone. Idealised C4MIP simulations, which can be used to separate the effects of CO2 and climate on the effective soil carbon turnover time, are very useful for assessing the role of false priming in models. Understanding these factors will be key to
predicting soil carbon changes over the next 100 years.

\citep{Varney2020}
Future changes in soil carbon depend on changes in litter and root inputs from plants and especially on reductions in the turnover time of soil carbon (taus) with warming. An approximation to the
latter term for the top one metre of soil (deltaCs,tau) can be diagnosed from projections made with the CMIP6 and CMIP5 Earth System Models (ESMs), and is found to span a large range even
at 2 degC of global warming (-196 +- 117 PgC). Here, we present a constraint on deltaCs,tau, which makes use of current heterotrophic respiration and the spatial variability of taus inferred from
observations. This spatial emergent constraint allows us to halve the uncertainty in deltaCs,tau at 2degC to -232 +- 52 PgC.

\citep{Wei2022}
Here, based on multiple global datasets and a traceability analysis, we diagnosed the uncertainty source of terrestrial carbon storage in 22 ESMs that participated in phases 5 and 6 of the Coupled Model Intercomparison Project (CMIP5 and CMIP6). The modeled global terrestrial carbon storage has
converged among ESMs from CMIP5 (1936.9 6 739.3 PgC) to CMIP6 (1774.4 6 439.0 PgC) but is persistently lower than the observation-based estimates (2285 6 669 PgC). By further decomposing terrestrial carbon storage into net primary production (NPP) and ecosystem carbon residence time (tE), we found that the decreased intermodel spread in land carbon storage primarily resulted from more accurate simulations on NPP among ESMs from CMIP5 to CMIP6. The persistent underestimation of land carbon storage was caused by the biased tE. In CMIP5 and CMIP6, the modeled tE was far shorter
than the observation-based estimates. The potential reasons for the biased tE could be the lack of or incomplete representation of nutrient limitation, vertical soil biogeochemistry, and the permafrost carbon cycle. Moreover, the modeled tE became the key driver for the intermodel spread in global land carbon storage in CMIP6.

The models in CMIP6 are adequate in capturing the near-present terrestrial productivity, while the modeled tE is lower than the dataset products. Thus, tE becomes the essential source for model ensemble bias and intermodel spread in terrestrial C storage.

\citep{Winkler2019} CMIP5

\citep{Zechlau2022}
. The spread of the sensitivity of tropical land carbon uptake to tropical
warming in an idealized simulation with a 1\% per year increase of atmospheric CO2 shows only a slight decrease in CMIP6 (−52 ± 35 GtC/K) compared to CMIP5 (−49 ± 40 GtC/K). For both model generations, the observed interannual variability in the growth rate of atmospheric CO2 yields a consistent emergent constraint on the sensitivity of tropical land carbon uptake with a constrained range of −37 ± 14 GtC/K for the combined ensemble (i.e., a reduction of ∼30% in the best estimate and 60% in the uncertainty range relative to the multimodel mean of the combined ensemble). A further emergent constraint is based on a relationship between CO2 fertilization and the historical increase in the CO2 seasonal cycle amplitude in high latitudes. However, this emergent constraint is not evident in CMIP6. This is in part because the historical increase in the amplitude of the CO2 seasonal cycle is more accurately simulated in CMIP6, such that the models are all now close to the observational constraint.

\citep{Zhu2022}
However, many uncertainties remain in projecting NPP responses, due to their complex processes and divergent model characteristics. This study estimated NPP responses to elevated CO2 and CO2-induced climate change using the Chinese Academy of Sciences Earth System Model version 2 (CAS-ESM2), as well as 22 CMIP6 models. Based on CMIP6 pre-industrial and abruptly quadrupled CO2 experiments, the analysis focused on a comparison of the CAS-ESM2 with the multi-model ensemble (MME), and on a detection of underlying causes of their differences. We found that all of the models showed an overall enhancement in NPP, and that CAS-ESM2 projected a slightly weaker NPP enhancement than MME. This weaker NPP enhancement was the net result of much weaker NPP enhancement over the tropics, and a little stronger NPP enhancement over northern high latitudes. We further report that these differences in NPP responses between the CAS-ESM2 and MME resulted from their different behaviors in simulating NPP trends with modeling time, and are attributed to their different projections of CO2-induced climatic anomalies and different climate sensitivities.
\bibliographystyle{elsarticle-harv}
\bibliography{sample}

\end{document}